% ============================================================
% CHAPITRE 8 — Sprint 5 : Module Réclamations
% ============================================================
\chapter{Sprint 5 --- Module Réclamations}

\section{Introduction}

Après le parcours de recrutement (offres, candidatures, entretiens et
embauche), ce sprint ajoute un \textbf{canal de feedback} transversal~:
les utilisateurs (candidats, RH, Admin) peuvent déposer une
\textbf{réclamation}, et le RH / Admin la traite jusqu'à clôture.

Conformément à l'architecture prévue au Sprint~0, les réclamations sont
implémentées dans un \textbf{microservice dédié} (\texttt{RECLAMATION},
port~8091), exposé via la Gateway sous le préfixe
\texttt{/api/reclamation/**}, avec la même sécurité JWT et le même appel
Feign vers le service \texttt{USER}.


\section{Objectifs du Sprint}

\begin{itemize}
  \item créer le microservice \texttt{RECLAMATION} (Spring Boot, Eureka,
        MySQL, JWT)~;
  \item permettre à un utilisateur authentifié de \textbf{créer} une
        réclamation (titre + description)~;
  \item permettre au candidat de \textbf{consulter} et
        \textbf{supprimer} ses réclamations encore ouvertes~;
  \item permettre au RH et à l'Admin de \textbf{lister} toutes les
        réclamations~;
  \item permettre au RH / Admin de \textbf{traiter} une réclamation
        (statut + réponse)~;
  \item permettre à l'Admin de \textbf{supprimer} une réclamation~;
  \item intégrer les écrans Angular côté candidat, RH et Admin.
\end{itemize}


\section{Spécification et Analyse des Besoins}

\subsection{Sprint Backlog}

\begin{table}[htbp]
  \centering
  \scriptsize
  \caption{Sprint Backlog du Sprint 5 (réclamations)}
  \label{tab:sprint5-backlog}
  \begin{tabular}{|p{2.1cm}|c|p{4.0cm}|c|p{3.1cm}|c|}
    \hline
    \textbf{Fonctionnalité} & \textbf{ID} & \textbf{Scénario utilisateur} &
    \textbf{ID Tâche} & \textbf{Tâches} & \textbf{Est.} \\
    \hline
    Microservice
    & US-22
    & En tant qu'équipe, je déploie le service réclamation autonome.
    & T1
    & Projet Spring Boot + Eureka + JWT
    & 1,5 j \\
    \hline
    Créer une réclamation
    & US-23
    & En tant qu'utilisateur, je dépose une réclamation.
    & T2
    & \texttt{POST /reclamations}
    & 1 j \\
    \cline{4-6}
    & & & T3 & Page candidat \texttt{/jobs/reclamations} & 1 j \\
    \hline
    Mes réclamations
    & US-24
    & En tant que candidat, je consulte et suis mes réclamations.
    & T4
    & \texttt{GET /reclamations/my}
    & 0,5 j \\
    \hline
    Traiter (RH/Admin)
    & US-25
    & En tant que RH/Admin, je change le statut et réponds.
    & T5
    & \texttt{PATCH .../status} + écrans RH/Admin
    & 1,5 j \\
    \hline
    Suppression
    & US-26
    & En tant qu'auteur (ouverte) ou Admin, je supprime une réclamation.
    & T6
    & \texttt{DELETE /reclamations/\{id\}}
    & 0,5 j \\
    \hline
    Gateway
    & US-27
    & Le front appelle le service via \texttt{/api/reclamation/**}.
    & T7
    & Route Gateway + menu Angular
    & 0,5 j \\
    \hline
  \end{tabular}
\end{table}


\subsection{Diagramme de cas d'utilisation}

\begin{figure}[H]
  \centering
  \reportfig{img/uc-sprint5.jpg}
  \caption{Diagramme de cas d'utilisation du Sprint 5}
  \label{fig:uc-sprint5}
\end{figure}


\section{Description textuelle des cas d'utilisation}

\subsection{Cas d'utilisation <<~Créer une réclamation~>>}

\begin{table}[htbp]
  \centering
  \small
  \caption{Description textuelle --- Créer une réclamation}
  \label{tab:uc-create-reclamation}
  \begin{tabular}{|p{3.0cm}|p{10.2cm}|}
    \hline
    \textbf{Titre} & Créer une réclamation \\
    \hline
    \textbf{Acteur} & Utilisateur authentifié (candidat, RH ou Admin) \\
    \hline
    \textbf{Précondition} & Compte actif, JWT valide. \\
    \hline
    \textbf{Postcondition} & Réclamation créée au statut \texttt{OPEN}. \\
    \hline
    \textbf{Scénario nominal}
    &
    \begin{enumerate}[leftmargin=*,itemsep=0.15em,topsep=0.2em]
      \item L'utilisateur ouvre \texttt{/jobs/reclamations} (candidat)
            ou l'espace RH/Admin.
      \item Il saisit un titre et une description.
      \item Le frontend envoie \texttt{POST /api/reclamation/reclamations}.
      \item Le service enregistre la réclamation liée à
            \texttt{createdByUserId}.
      \item La réclamation apparaît dans \texttt{Mes réclamations}.
    \end{enumerate}
    \\
    \hline
  \end{tabular}
\end{table}


\subsection{Cas d'utilisation <<~Traiter une réclamation~>>}

\begin{table}[htbp]
  \centering
  \small
  \caption{Description textuelle --- Traiter une réclamation}
  \label{tab:uc-process-reclamation}
  \begin{tabular}{|p{3.0cm}|p{10.2cm}|}
    \hline
    \textbf{Titre} & Traiter une réclamation \\
    \hline
    \textbf{Acteur} & RH ou Admin \\
    \hline
    \textbf{Précondition} & Au moins une réclamation existante. \\
    \hline
    \textbf{Postcondition} & Statut mis à jour et réponse RH éventuellement renseignée. \\
    \hline
    \textbf{Scénario nominal}
    &
    \begin{enumerate}[leftmargin=*,itemsep=0.15em,topsep=0.2em]
      \item Le RH ouvre \texttt{/rh/reclamations} (ou Admin
            \texttt{/admin/reclamations}).
      \item Il consulte la liste (auteur, titre, statut, date).
      \item Il clique sur \textbf{Traiter}, choisit un statut
            (\texttt{OPEN}, \texttt{IN\_PROGRESS}, \texttt{RESOLVED},
            \texttt{CLOSED}) et saisit une réponse.
      \item \texttt{PATCH /api/reclamation/reclamations/\{id\}/status}
            met à jour la fiche et mémorise le traiteur
            (\texttt{handledByUserId}).
    \end{enumerate}
    \\
    \hline
  \end{tabular}
\end{table}


\subsection{Cas d'utilisation <<~Consulter mes réclamations~>>}

\begin{table}[htbp]
  \centering
  \small
  \caption{Description textuelle --- Consulter mes réclamations}
  \label{tab:uc-my-reclamations}
  \begin{tabular}{|p{3.0cm}|p{10.2cm}|}
    \hline
    \textbf{Titre} & Consulter mes réclamations \\
    \hline
    \textbf{Acteur} & Candidat \\
    \hline
    \textbf{Précondition} & Candidat authentifié. \\
    \hline
    \textbf{Postcondition} & Liste personnelle affichée avec réponses RH. \\
    \hline
    \textbf{Scénario nominal}
    &
    \begin{enumerate}[leftmargin=*,itemsep=0.15em,topsep=0.2em]
      \item Le candidat ouvre \texttt{/jobs/reclamations}.
      \item \texttt{GET /api/reclamation/reclamations/my} retourne ses
            réclamations triées par date.
      \item Il peut supprimer une réclamation encore \texttt{OPEN}.
    \end{enumerate}
    \\
    \hline
  \end{tabular}
\end{table}


\subsection{Autres cas d'utilisation}

\begin{itemize}
  \item \textbf{Modifier une réclamation ouverte}~: l'auteur peut mettre
        à jour titre / description tant que le statut est \texttt{OPEN}.
  \item \textbf{Supprimer (Admin)}~: suppression définitive depuis
        \texttt{/admin/reclamations}.
  \item \textbf{Consulter le détail}~:
        \texttt{GET /reclamations/\{id\}} réservé au propriétaire, au RH
        ou à l'Admin.
\end{itemize}


\section{Conception}

\subsection{Architecture du microservice}

Le service \texttt{RECLAMATION} suit le même modèle que
\texttt{RECRUITMENT}~:
\begin{itemize}
  \item validation JWT (secret partagé) et résolution de l'utilisateur via
        Feign \texttt{USER}~;
  \item base MySQL dédiée \texttt{reclamations}~;
  \item enregistrement Eureka et routage Gateway
        \texttt{/api/reclamation/**} $\rightarrow$ \texttt{lb://RECLAMATION}.
\end{itemize}

\subsection{Diagramme de classes}

Le domaine est volontairement simple~: une entité \textbf{Reclamation}
avec statut, auteur, traiteur et réponse RH.

\begin{figure}[H]
  \centering
  \reportfig{img/class-sprint5.jpg}
  \caption{Diagramme de classes du Sprint 5}
  \label{fig:class-sprint5}
\end{figure}


\subsection{Diagrammes de séquence}

\subsubsection{Diagramme de séquence <<~Créer une réclamation~>>}

\begin{figure}[H]
  \centering
  \reportfig{img/seq-create-reclamation.jpg}
  \caption{Diagramme de séquence --- Créer une réclamation}
  \label{fig:seq-create-reclamation}
\end{figure}

\subsubsection{Diagramme de séquence <<~Traiter une réclamation~>>}

\begin{figure}[H]
  \centering
  \reportfig{img/seq-process-reclamation.jpg}
  \caption{Diagramme de séquence --- Traiter une réclamation}
  \label{fig:seq-process-reclamation}
\end{figure}


\section{Réalisation}

\subsection{API REST}

\begin{table}[htbp]
  \centering
  \small
  \caption{Endpoints principaux du microservice Réclamation}
  \label{tab:api-reclamation}
  \begin{tabular}{|l|p{6.2cm}|p{3.8cm}|}
    \hline
    \textbf{Méthode} & \textbf{Endpoint} & \textbf{Accès} \\
    \hline
    POST & \texttt{/api/reclamation/reclamations} & Authentifié \\
    \hline
    GET & \texttt{/api/reclamation/reclamations/my} & Authentifié \\
    \hline
    GET & \texttt{/api/reclamation/reclamations} & RH / Admin \\
    \hline
    GET & \texttt{/api/reclamation/reclamations/\{id\}} & Propriétaire / RH / Admin \\
    \hline
    PUT & \texttt{/api/reclamation/reclamations/\{id\}} & Propriétaire (OPEN) / RH / Admin \\
    \hline
    PATCH & \texttt{/api/reclamation/reclamations/\{id\}/status} & RH / Admin \\
    \hline
    DELETE & \texttt{/api/reclamation/reclamations/\{id\}} & Propriétaire (OPEN) / Admin \\
    \hline
  \end{tabular}
\end{table}

Les statuts gérés sont~: \texttt{OPEN}, \texttt{IN\_PROGRESS},
\texttt{RESOLVED}, \texttt{CLOSED}.


\subsection{Interface candidat}

Le candidat crée et suit ses réclamations depuis
\texttt{/jobs/reclamations}.

\begin{figure}[H]
  \centering
  \reportfig{img/ui-my-reclamations.png}
  \caption{Interface candidat --- mes réclamations}
  \label{fig:ui-my-reclamations}
\end{figure}


\subsection{Interface RH}

Le RH traite les réclamations (statut + réponse) depuis
\texttt{/rh/reclamations}.

\begin{figure}[H]
  \centering
  \reportfig{img/ui-rh-reclamations.png}
  \caption{Interface RH --- traitement des réclamations}
  \label{fig:ui-rh-reclamations}
\end{figure}


\subsection{Interface Admin}

L'Admin dispose de la même liste, avec suppression définitive.

\begin{figure}[H]
  \centering
  \reportfig{img/ui-admin-reclamations.png}
  \caption{Interface Admin --- gestion des réclamations}
  \label{fig:ui-admin-reclamations}
\end{figure}


\section{Conclusion}

Le Sprint~5 a livré un microservice \textbf{Réclamation} autonome,
intégré à la Gateway et aux espaces candidat, RH et Admin. Le cycle de
vie est simple et opérationnel~: dépôt $\rightarrow$ traitement
$\rightarrow$ résolution / clôture.

Le chapitre suivant présentera l'\textbf{industrialisation DevOps}
(Docker, Jenkins et SonarQube) de la plateforme DAAM.

\clearpage
